\documentclass[12pt]{ctexart}

% ====== 基本设置 ======
\usepackage[margin=1in]{geometry}
\usepackage{booktabs,longtable,tabularx,array}
\usepackage{hyperref}
\usepackage{graphicx}
\usepackage{enumitem}
\usepackage{float}
\usepackage{xcolor}

\hypersetup{
  colorlinks=true,
  linkcolor=blue,
  urlcolor=blue,
  citecolor=blue
}

\pagestyle{plain}
\pagenumbering{arabic}

\newcolumntype{L}{>{\raggedright\arraybackslash}X}


\newcommand{\ph}[1]{\textit{【#1】}}

\begin{document}

\section{本文使用AI工具名称及版本对应}

\begin{table}[!htbp]
    \caption{本文使用AI工具名称及版本对应}\label{tab:001} \centering
    \begin{tabular}{ccccc}
        \toprule[1.5pt]
    \textbf{工具名称} & \textbf{版本号} \\
    \midrule[1pt]
    Deepseek & Deepseek-R1-0528 \\
    ChatGPT & GPT-5 \\
    Claude & 4.0 Sonnet \\
    豆包 & 1.6 \\
    文心一言 & 4.5 Turbo \\
    Kimi & K2 \\
    \bottomrule[0.5pt]
    \end{tabular}
\end{table}
\ph{注：用了多少个AI就在这里标注多少个，下面的附录中每个AI单作一章。}

\newpage

\section{GPT-5 使用详情}
\subsection{具体使用目的及环节}

\subsubsection{文本润色}
\paragraph{使用环节一}
章节：\ph{模型评价与推广——模型的优点} \\
目的：\ph{用于文本的润色，使其按照要求呈现形式更加xx化}

\paragraph{使用环节二}
章节：\ph{在此填写章节} \\
目的：\ph{在此概述润色目标与改写侧重点}

\subsubsection{代码辅助}
\paragraph{使用环节一}
章节：\ph{模型建立与求解——问题一模型建立与求解} \\
目的：\ph{提供了问题一中xx部分的代码，采纳了xx部分用于辅助解题}

\paragraph{使用环节二}
章节：\ph{在此填写章节} \\
目的：\ph{在此描述辅助编码内容与边界}

\subsubsection{图表生成}
\paragraph{使用环节一}
章节：\ph{模型建立与求解——问题一模型建立与求解} \\
目的：\ph{提供了问题一中xx部分的图/表，用于完成xx的可视化/分析}

\paragraph{使用环节二}
章节：\ph{在此填写章节} \\
目的：\ph{在此描述图表用途与产出格式}

\subsubsection{思路辅助}
\paragraph{使用环节一}
章节：\ph{模型建立与求解——问题一模型建立与求解} \\
目的：\ph{提供了问题一中xx部分的思路，用于辅助问题一中的数据处理/模型建立/xxx}

\paragraph{使用环节二}
章节：\ph{在此填写章节} \\
目的：\ph{在此描述思路启发与采用范围}

\subsection{关键交互记录及采纳情况}
\subsubsection{交互记录一}
\textbf{交互场景}：\ph{见 2.1.1} \\[4pt]
\textbf{交互记录截图}：\ph{插入截图}

\paragraph{提示词}
\begin{flushleft}
\ph{在此粘贴提示词文本}
\end{flushleft}

\paragraph{回复内容}
\begin{flushleft}
\ph{在此粘贴模型回复要点或全文}
\end{flushleft}

\paragraph{采纳部分}
\begin{itemize}[leftmargin=1.5em]
  \item \ph{列出被直接采纳的段落/算法/图表}
\end{itemize}

\paragraph{修改部分}
\begin{itemize}[leftmargin=1.5em]
  \item \ph{列出在采纳基础上的修订点及理由}
\end{itemize}

\paragraph{最终呈现}
\begin{flushleft}
\ph{说明最终落地内容、位置与版本（例如章节、页码、图号等）}
\end{flushleft}

\newpage

\section{Deepseek-R1-0528 使用详情}
\subsection{具体使用目的及环节}

\subsubsection{文本润色}
\paragraph{使用环节一}
章节：\ph{模型评价与推广——模型的优点} \\
目的：\ph{用于文本的润色，使其按照要求呈现形式更加xx化}

\paragraph{使用环节二}
章节：\ph{在此填写章节} \\
目的：\ph{在此概述润色目标与改写侧重点}

\subsubsection{代码辅助}
\paragraph{使用环节一}
章节：\ph{模型建立与求解——问题一模型建立与求解} \\
目的：\ph{提供了问题一中xx部分的代码，采纳了xx部分用于辅助解题}

\paragraph{使用环节二}
章节：\ph{在此填写章节} \\
目的：\ph{在此描述辅助编码内容与边界}

\subsubsection{图表生成}
\paragraph{使用环节一}
章节：\ph{模型建立与求解——问题一模型建立与求解} \\
目的：\ph{提供了问题一中xx部分的图/表，用于完成xx的可视化/分析}

\paragraph{使用环节二}
章节：\ph{在此填写章节} \\
目的：\ph{在此描述图表用途与产出格式}

\subsubsection{思路辅助}
\paragraph{使用环节一}
章节：\ph{模型建立与求解——问题一模型建立与求解} \\
目的：\ph{提供了问题一中xx部分的思路，用于辅助问题一中的数据处理/模型建立/xxx}

\paragraph{使用环节二}
章节：\ph{在此填写章节} \\
目的：\ph{在此描述思路启发与采用范围}

\subsection{关键交互记录及采纳情况}
\subsubsection{交互记录一}
\textbf{交互场景}：\ph{见 3.1.1} \\[4pt]
\textbf{交互记录截图}：\ph{插入截图}

\paragraph{提示词}
\begin{flushleft}
\ph{在此粘贴提示词文本}
\end{flushleft}

\paragraph{回复内容}
\begin{flushleft}
\ph{在此粘贴模型回复要点或全文}
\end{flushleft}

\paragraph{采纳部分}
\begin{itemize}[leftmargin=1.5em]
  \item \ph{列出被直接采纳的段落/算法/图表}
\end{itemize}

\paragraph{修改部分}
\begin{itemize}[leftmargin=1.5em]
  \item \ph{列出在采纳基础上的修订点及理由}
\end{itemize}

\paragraph{最终呈现}
\begin{flushleft}
\ph{说明最终落地内容、位置与版本}
\end{flushleft}

\newpage

\section{Claude-4.0 Sonnet 使用详情}
\subsection{具体使用目的及环节}

\subsubsection{文本润色}
\paragraph{使用环节一}
章节：\ph{模型评价与推广——模型的优点} \\
目的：\ph{用于文本的润色，使其按照要求呈现形式更加xx化}

\paragraph{使用环节二}
章节：\ph{在此填写章节} \\
目的：\ph{在此概述润色目标与改写侧重点}

\subsubsection{代码辅助}
\paragraph{使用环节一}
章节：\ph{模型建立与求解——问题一模型建立与求解} \\
目的：\ph{提供了问题一中xx部分的代码，采纳了xx部分用于辅助解题}

\paragraph{使用环节二}
章节：\ph{在此填写章节} \\
目的：\ph{在此描述辅助编码内容与边界}

\subsubsection{图表生成}
\paragraph{使用环节一}
章节：\ph{模型建立与求解——问题一模型建立与求解} \\
目的：\ph{提供了问题一中xx部分的图/表，用于完成xx的可视化/分析}

\paragraph{使用环节二}
章节：\ph{在此填写章节} \\
目的：\ph{在此描述图表用途与产出格式}

\subsubsection{思路辅助}
\paragraph{使用环节一}
章节：\ph{模型建立与求解——问题一模型建立与求解} \\
目的：\ph{提供了问题一中xx部分的思路，用于辅助问题一中的数据处理/模型建立/xxx}

\paragraph{使用环节二}
章节：\ph{在此填写章节} \\
目的：\ph{在此描述思路启发与采用范围}

\subsection{关键交互记录及采纳情况}
\subsubsection{交互记录一}
\textbf{交互场景}：\ph{见 4.1.1} \\[4pt]
\textbf{交互记录截图}：\ph{插入截图}

\paragraph{提示词}
\begin{flushleft}
\ph{在此粘贴提示词文本}
\end{flushleft}

\paragraph{回复内容}
\begin{flushleft}
\ph{在此粘贴模型回复要点或全文}
\end{flushleft}

\paragraph{采纳部分}
\begin{itemize}[leftmargin=1.5em]
  \item \ph{列出被直接采纳的段落/算法/图表}
\end{itemize}

\paragraph{修改部分}
\begin{itemize}[leftmargin=1.5em]
  \item \ph{列出在采纳基础上的修订点及理由}
\end{itemize}

\paragraph{最终呈现}
\begin{flushleft}
\ph{说明最终落地内容、位置与版本}
\end{flushleft}

\end{document}
